\documentclass{report}
\usepackage{amsmath,amssymb}
\setlength{\parindent}{0mm}
\setlength{\parskip}{1em}
\begin{document}
\begin{center}
\rule{6in}{1pt} \
{\large Karl Rupp \\
{\bf A Performance Comparison of Algebraic Multigrid Preconditioners on GPUs and MIC}}

Institute for Microelectronics \\ Vienna University of Technology \\ Gu\ss{}hausstra\ss{}e 27-29/E360 \\ A-1040 Wien \\ Austria \\ Europe
\\
{\tt rupp@iue.tuwien.ac.at}\\
Ansgar J\"ungel\\
Tibor Grasser\end{center}

Algebraic multigrid (AMG) preconditioners for accelerators such as
graphics processing units (GPUs) and Intel's many-integrated core (MIC)
architecture typically require a careful, problem-dependent trade-off
between efficient hardware use, robustness, and convergence rate in order
to minimize time-to-solution. Several variants of AMG with fine-grained
parallelism have been proposed recently, but a comparison across
different hardware architectures is difficult since the proposed
approaches are mostly focused on a single architecture. To address this
deficiency, we derived implementations of recently proposed AMG variants
in CUDA, OpenCL, and OpenMP and run extensive benchmarks. Our performance
results for GPUs from AMD and NVIDIA as well as for Intel's Xeon Phi
reveal the sweet spots of each accelerator architecture, helping
practitioners to select the best AMG variant and hardware for a given
problem.


\end{document}
